\section{The Selmer sign and actual ordinary adjoint points}
\label{sec:kt-host-sign}

This section uses the same pinned-source convention as
KT-ONECHAR-CLOSED.  The new local argument below is a mathematical
proof, not a claimed Lean compilation.  Put
$E=\mathbf Q(\sqrt{-3})$, $\eta=\chi_{-3}$, and
$\Psi=\varepsilon_\ell^e$, with the arithmetic cyclotomic
normalization.  Assume $\ell>3$, $e\geq2$ is even, and
$(\ell-1)\nmid2e$.  These are the relevant rational trivial-character
\texttt{GD} conditions.  Write $J$ for the anti-diagonal matrix with
all three anti-diagonal entries equal to one.  The source uses
\[
 \theta(g)=Jg^{-t}J,\qquad
 GG_3=\mathrm{GL}_3\rtimes\{\pm1\},\qquad
 r_0(g)=(\operatorname{diag}(\Psi(g),1,\Psi(g)^{-1}),\eta(g)).
                                                               \tag{HS.1}
\]
Here $-1$ acts by $\theta$.  Matrix units in this section use
indices $1,2,3$; the Lean files use $0,1,2$.

\begin{lemma}[KT-HOST-TANGENT-SIGN]
The weight $-1$ space for the adjoint action of $r_0$ splits as
\[
 \langle E_{21},E_{32}\rangle
  =\langle E_{21}-E_{32}\rangle
       \oplus\langle E_{21}+E_{32}\rangle,
                                                               \tag{HS.2}
\]
with global characters $\Psi^{-1}$ and $\eta\Psi^{-1}$,
respectively.  The rational strict Selmer injection used by
\texttt{ShaInjGen} takes its values in the first line.  The lower
root of the normalized principal adjoint embedding takes its values
in the second line.
\end{lemma}
\begin{proof}
The differential of the involution is
$\theta_{\rm Lie}(X)=-JX^tJ$, so it interchanges
$E_{21}$ and $-E_{32}$.  Conjugation by
$\operatorname{diag}(z,1,z^{-1})$ multiplies both matrix units
by $z^{-1}$.  This proves the two characters in (HS.2).
In particular complex conjugation acts by $+1$ and $-1$ on
the two lines, because $\Psi(c)=1$ and $\eta(c)=-1$.

The literal definition \texttt{shaInj\_mat} is
$t\mapsto t(E_{21}-E_{32})$; the following theorem
\texttt{shaInj\_ad} proves its equivariance with character
$\Psi^{-1}$ for every global group element.
For the second assertion, use the ordered basis
\[
 E=\begin{pmatrix}0&1\\0&0\end{pmatrix},\qquad
 H=\begin{pmatrix}1&0\\0&-1\end{pmatrix},\qquad
 F=\begin{pmatrix}0&0\\1&0\end{pmatrix}
\]
of $\mathfrak{sl}_2$ and replace $H$ by $sH$, where
$s^2=-1/2$.  Then
\[
 \operatorname{ad}(F)=-s^{-1}E_{21}+2sE_{32}
                  =-s^{-1}(E_{21}+E_{32}).                    \tag{HS.3}
\]
The preserved trace form in this basis is
\[
 K=\begin{pmatrix}0&0&1\\0&-1&0\\1&0&0\end{pmatrix}.
                                                               \tag{HS.4}
\]
The construction below identifies this basis with the involution
in (HS.1), proving the asserted normalization.
\end{proof}

The distinction in (HS.2) is a distinction between tangent
directions.  To decide whether it excludes a characteristic-zero
point, one must check the actual point conditions.  In particular,
\texttt{TypeI} means irreducibility on $G_E$; it does not require
that a point arise by exponentiating the first line in (HS.2).
The next lemma addresses the local condition that is missing from
an argument using residual semisimplification alone.

\begin{lemma}[KT-ORDINARY-ON-BASE-COMPONENT]
Let $O$ be the ring of integers of any finite extension of
$\mathbf Q_\ell$,
and let $D_{\rm loc}$ be the source's local data with unramified
$\eta$, trivial $\tau$, and $\Psi|_I=\varepsilon_\ell^e$.
Let $U$ be its unrestricted framed universal ring.  Assume the
carried local component theorem
\[
 \mathrm{primesAt0}(D_{\rm loc},U)=\{I_t\}.                  \tag{HS.5}
\]
Let $A$ be a complete Noetherian local $O$-algebra with the same
residue field.  Suppose an integral lift of $r_0\bmod\mathfrak m_O$
to $GG_3(A)$ is conjugate, by a matrix congruent to the identity,
to a lift valued in $B\rtimes\{\pm1\}$, where $B$ is the standard
upper Borel, and its diagonal character on inertia is exactly
$\operatorname{diag}(\varepsilon_\ell^e,1,
\varepsilon_\ell^{-e})$.  Then its map $U.R\to A$ kills $I_t$.
This holds whether the quadratic local extension is split or
unramified quadratic.
\end{lemma}

\begin{proof}
\paragraph{A smooth integral family with the prescribed inertia.}
First consider the framed functor of lifts valued in
$B\rtimes\{\pm1\}$, with the projection $\eta$ fixed and the
specified torus character on inertia fixed.  It is prorepresented
by a complete Noetherian local $O$-algebra: this is the usual
framed lifting construction, imposing the Borel and inertial
matrix equations.  Its finite tangent space follows from local
Galois cohomological finiteness.  The unipotent radical $N$,
its Lie algebra $\mathfrak n$, the Borel Lie algebra
$\mathfrak b$, and $\mathfrak t=\mathfrak b/\mathfrak n$
are all stable under $\theta$.

Here is the full lifting argument, including the unramified
quadratic case.  For a small surjection $A'\to A$ with kernel
$J$, the obstruction to lifting a homomorphism with fixed
projection $\eta$ is the usual two-cocycle in
$H^2(G_{\mathbf Q_\ell},\mathfrak b_{\bar r_0})\otimes J$.
The fixed outer action merely makes this the indicated twisted
adjoint module; it does not introduce an additional Frobenius
descent condition.  On inertia the composition factors of
$\mathfrak b$ have characters
\[
                   1,\quad\bar\varepsilon_\ell^e,
                         \quad\bar\varepsilon_\ell^{2e}.
                                                               \tag{HS.6}
\]
Every exponent is even.  The Tate dual of each factor has an
odd exponent, which is nonzero modulo the even number
$\ell-1$.  Local Tate duality therefore gives
\[
 H^2(G_{\mathbf Q_\ell},\mathfrak b_{\bar r_0})
   =H^2(G_{\mathbf Q_\ell},\mathfrak n_{\bar r_0})=0.
                                                               \tag{HS.7}
\]
This proves that a Borel lift exists over $A'$.  Its torus
restriction may differ from the required inertia character by
an element of $H^1(I,\mathfrak t)^{G/I}\otimes J$.
The correction map
\[
 H^1(G,\mathfrak b)\longrightarrow
 H^1(G,\mathfrak t)\longrightarrow
 H^1(I,\mathfrak t)^{G/I}                                  \tag{HS.8}
\]
is surjective: the first arrow is surjective by (HS.7), and
the second by $\operatorname{cd}_\ell(G/I)=1$.
Changing the lift by a one-cocycle consequently corrects its
inertia without changing the given reduction.  Here inertia acts
trivially on $\mathfrak t$, so equality of its cohomology classes
is equality of the actual torus correction cocycles; there is
no unaccounted inertia coboundary.  Thus the functor
is formally smooth and its representing ring $R_B$ is a power
series ring over $O$, in particular a domain.  No unramifiedness
of the coefficient DVR is used: (HS.6)--(HS.8) are statements
over its finite residue field, and formal smoothness is relative
to the given $O$.

For split $\eta$ this is precisely the lifting argument in
Patrikis, \emph{Deformations of Galois representations and
exceptional monodromy}, \S4.1, Proposition~4.4, preprint
pp.~12--13.  Its hypotheses \textnormal{(REG)} and
\textnormal{(REG*)} hold here: the negative-root inertia
characters are $\bar\varepsilon_\ell^{-e}$ and
$\bar\varepsilon_\ell^{-2e}$, nontrivial by
$(\ell-1)\nmid2e$, and their cyclotomic twists have odd
nonzero exponents.  The preceding obstruction calculation
provides the necessary adaptation for nonsplit unramified
$\eta$ directly.  In particular we have not substituted a
naive fixed-point condition for descent of a homomorphism.

Adjoin nine formal variables to $R_B$ and conjugate its
universal lift by the matrix $1+(X_{ij})$.  The resulting
domain $R'=R_B[[X_{ij}]]$ is a family containing every
congruently conjugate lift in the lemma.  It also contains
the exact base representation $r_0$.

\paragraph{The saturation and the chosen component.}
Recall the literal definitions in the source:
$I_Y=(I_{\rm ord}:\delta^\infty)$,
$I_{\rm cr}=\sqrt{I_Y}$, and $I_t$ is the intersection of
the minimal primes of $I_{\rm cr}$ lying under the base-point
kernel.  Here $I_{\rm ord}$ is generated by the prescribed
inertia characteristic polynomials and the entries of
\[
 (r(s_0)-\Psi(s_0))
 (r(s_1)-1)
 (r(s_2)-\Psi(s_2)^{-1}),\qquad s_i\in I.                   \tag{HS.9}
\]
Every expression in (HS.9) vanishes on an upper triangular
lift in the indicated order: the rightmost factor maps the
three-dimensional space to the second flag step, the next
maps that step to the first, and the leftmost kills the first.
The characteristic polynomial equations also vanish.  These
statements are preserved under conjugation, so the map
$\phi:U.R\to R'$ kills $I_{\rm ord}$.

The image of the saturation element $\delta$ in $R'$ is
nonzero.  To check this without assuming it is nonzero at
the base, specialize to the diagonal family
$r_0(\eta,\Psi u_T)$, where $u_T$ is unramified and sends
a chosen Frobenius to $1+T$.  If the source's chosen
$\mathrm{fr}$ has Frobenius exponent $m\geq1$, put
$y=\Psi(\mathrm{fr})(1+T)^m$ and
$v=\varepsilon_\ell(\mathrm{fr})$.  Writing $c_i$ for the
source's nonzero products of distinct inertia eigenvalue
differences, direct substitution gives
\[
 \delta=\ell\,c_0c_1^2c_2\,y^{-1}(y-v)^2\ne0
                                  \quad\text{in }O[[T]].     \tag{HS.10}
\]
The $c_i$ are nonzero because the chosen inertia element has
cyclotomic character of infinite order.  The function $y$ is
nonconstant and the coefficient ring has characteristic zero.
Formula (HS.10) is also the computation used in the archived
\texttt{del\_unr\_twist} proof.  Since $R'$ is a domain,
$\phi$ kills $I_Y$ and hence $I_{\rm cr}$.

Set $P=\ker\phi$.  It is prime.  The base specialization
$R'\to O$ shows $P\subseteq\ker x_0$, where $x_0$ is the
unique framed base-point map.  Choose a minimal prime $Q$ of
$I_{\rm cr}$ contained in $P$.  Such a prime exists since
$U.R$ is Noetherian.  Then $Q\subseteq\ker x_0$, so $Q$
belongs to \texttt{primesAt0}.  By (HS.5), $Q=I_t$.
Thus $I_t\subseteq P$.  Specializing the family to the lift
in the lemma proves the assertion.
\end{proof}

\begin{theorem}[KT-ADJOINT-ACTUAL-HOST-POINT]
Let $f\in S_{e+1}(\Gamma_0(3),\eta)$ be a non-CM primitive
eigenform and let $\lambda\mid\ell$ be a coefficient prime.
Suppose $a_\ell(f)$ is a $\lambda$-adic unit and
\[
 \bar r_f^{\,\mathrm{ss}}
       =1\oplus\eta\bar\varepsilon_\ell^e.                 \tag{HS.11}
\]
Keep the assumptions at the beginning of this section and
the carried local component theorem (HS.5).  For any of the
global data $D$, an element $g_1$ satisfying $D.\mathrm{Marked}(g_1)$,
local universal ring
$U$, and global ring $S=D.\mathrm{Ring}(U,g_1)$ used in the
rational trivial-character \texttt{OneCharU\_s} construction,
there is a point
\[
                  z:S.R\longrightarrow\overline{\mathbf Q}_\ell
                                                               \tag{HS.12}
\]
satisfying the literal predicates \texttt{IsLocPt} and
\texttt{TwoPt}.  Its representation on $G_E$ is isomorphic
to $\operatorname{Ad}(r_f)|_{G_E}$ and is of \texttt{TypeI}.
Consequently it also satisfies the corresponding
$D.\mathrm{IsPt}$ conditions.  This conclusion concerns
actual points of the prescribed source ring, not only
membership in the ambient admissible class.
\end{theorem}

\begin{proof}
\paragraph{A split residual lattice, with its polarization.}
Use the cohomological dual $r_f^\vee$.  Its residual
semisimplification is $1\oplus\eta\bar\Psi^{-1}$; its
adjoint is isomorphic to that of $r_f$.  Choose a stable
lattice over a finite coefficient DVR.  Because two is a
unit and complex conjugation has distinct eigenvalues
$1,-1$, the lattice has a basis in which
$r_f^\vee(c)=\operatorname{diag}(1,-1)$.  The residual
representation is triangular in this basis: an invariant
residual line must be one of the two distinct $c$-eigenlines.
Thus one of the two off-diagonal matrix-entry ideals is
contained in the maximal ideal.

Enlarge the coefficient DVR by a square root of a uniformizer
$\pi$.  If the upper off-diagonal entries are arbitrary
integral and the lower ones lie in $(\pi)$, replace the
second basis vector by $\sqrt\pi$ times itself.  The new
upper entries are multiplied by $\sqrt\pi$ and the new
lower entries divided by it.  Both are now in the maximal
ideal.  In the opposite triangular orientation use the
opposite scaling.  This preserves the diagonal action of
$c$ and yields an integral representation whose reduction
is exactly diagonal.  This argument does not assume that
a split residual lattice existed over the original DVR.

Enlarge coefficients once more to contain a unit $s$ with
$s^2=-1/2$, and form the adjoint in the ordered basis
$(E,sH,F)$ of (HS.3).  Denote this representation by $A(g)$.
It preserves the perfect form $K$ of (HS.4), and
\[
 A(c)=D_c:=\operatorname{diag}(-1,1,-1),\qquad
 JK=-D_c,\qquad\theta(A(g))=D_cA(g)D_c.                    \tag{HS.13}
\]
If $j(g)\in\{0,1\}$ is defined by
$\eta(g)=(-1)^{j(g)}$, set
\[
                 r(g)=(A(g)D_c^{j(g)},\eta(g)).              \tag{HS.14}
\]
Equation (HS.13) proves directly that (HS.14) is a
homomorphism to $GG_3$.  It has $r(c).\mathrm{mat}=1$.
The residual adjoint matrix is
$\operatorname{diag}(\eta\bar\Psi,1,
\eta\bar\Psi^{-1})$, and multiplication by $D_c^{j(g)}$
cancels the two $\eta$ factors.  Hence the residual matrix
of $r$ is exactly $r_0$, entry by entry.

Let $O$ be the coefficient DVR of $D$, enlarge $O'$ if necessary
so that it contains $O$, and take
$A_0=O+\mathfrak m_{O'}$ inside the final coefficient DVR $O'$.
Every entry of $r$ and of its inverse lies in $A_0$, because
its reduction is in the residue field of $O$.
This order is finite over $O$, local with maximal ideal
$\mathfrak m_{O'}$, and complete; its residue field is exactly
that of $O$.  Indeed it is the inverse image of that residue
subfield, and its maximal-ideal topology agrees with its
$\ell$-adic topology.  Thus it is an admissible coefficient
algebra for the source's \texttt{IsCO} universal property.
The unit $s$ itself need not lie in $A_0$; the adjoint
matrices do, which is all that is required.

\paragraph{Compatibility of the ordinary flag and the involution.}
The ordinary local theorem gives for $r_f^\vee$ an actual
upper triangular flag ordered as an unramified character
followed by a character with inertia
$\varepsilon_\ell^{-e}$.  Its integral saturated line
reduces to the first standard line: the two residual
inertia characters are distinct by
$(\ell-1)\nmid2e$.  A matrix $q\equiv1\pmod{\mathfrak m_{O'}}$
therefore makes $q r_f^\vee q^{-1}$ upper triangular
on $G_{\mathbf Q_\ell}$.  Its adjoint $h=\operatorname{Ad}(q)$
in the basis above also satisfies $h\equiv1$, lies in
$\mathrm{GL}_3(A_0)$, preserves $K$, and obeys
$\theta(h)=D_chD_c$.  Consequently
\[
 (\operatorname{conj}_h r)(g)
       =(hA(g)h^{-1}D_c^{j(g)},\eta(g))                     \tag{HS.15}
\]
is valued in $B\rtimes\{\pm1\}$ locally at $\ell$.
Since $\eta$ is unramified there, its diagonal inertia
character is exactly
$(\varepsilon_\ell^e,1,\varepsilon_\ell^{-e})$.
This proves compatibility with the polarization, including
the inert quadratic case.  KT-ORDINARY-ON-BASE-COMPONENT
puts the local representation on $I_t$.

\paragraph{The global normalization and universal point.}
The representation $r$ kills the exact subgroup
$D.\mathrm{Kill}$: away from $\ell$, its restriction to
$G_E$ is unramified by the already checked local conductor
calculation for the adjoint lift, and membership in
$D.\mathrm{Kill}$ forces $\eta(g)=1$.  Continuity and all
integrality conditions hold over $A_0$.

Apply the source's \texttt{exists\_normalised} theorem
over $A_0$ to the residual marked data.  It supplies a
matrix $v\equiv1$ such that $r'=\operatorname{conj}_v r$
is \texttt{Normalised} at $(c,g_1)$.  The local condition
proved above is preserved: locally $r'$ is still conjugate
by a matrix congruent to one to (HS.15), so the same local
lemma applies.  Killing $D.\mathrm{Kill}$ is also unchanged.
Thus $r'$ is an integral lift satisfying precisely
\texttt{NormCond(BoxCond)}.

The defining universal property of $S$ gives a local
$O$-algebra map $S.R\to A_0$ whose universal representation
is $r'$.  Composing with $A_0\hookrightarrow
\overline{\mathbf Q}_\ell$ gives (HS.12).
Every element of the maximal ideal of $S.R$ maps into
$\mathfrak m_{O'}$, so \texttt{IsLocPt} holds.
The representation on $G_E$ is isomorphic to
$\operatorname{Ad}(r_f)|_{G_E}$.  Its absolute
irreducibility is the all-prime non-CM adjoint
irreducibility proved in KT-ORDINARY-SYM2-FAMILY.
It is therefore of \texttt{TypeI} and can be neither
\texttt{Upper} nor \texttt{Lower}, proving \texttt{TwoPt}.
The source theorem \texttt{isPt\_push} now gives
$D.\mathrm{IsPt}$ as well; alternatively each clause was
verified directly in the construction.
\end{proof}

\begin{corollary}[KT-ADJOINT-HOST-OBSTRUCTION]
Every ordinary symmetric-square family in
KT-ORDINARY-SYM2-GD gives actual \texttt{TypeI},
\texttt{TwoPt} points as in KT-ADJOINT-ACTUAL-HOST-POINT.
In particular, under the same carried arithmetic inputs,
for the certified prime $\ell=697446897311$ and
\[
                   e=26+697446897310\,t,\qquad t\geq0,
\]
these points coexist with $K_{2e}(\mathbf Z)\{\ell\}=0$.
Their trace difference from the Eisenstein triple is a
nonunit at every element of $G_E$.  Any finite set satisfying
the literal \texttt{CoversPt} condition contains a
representative of their isomorphism class.
\end{corollary}
\begin{proof}
KT-ORDINARY-SYM2-FAMILY supplies the ordinary forms and
(HS.11); KT-ORDINARY-SYM2-GD supplies
$(\ell-1)\nmid2e$.  Apply the preceding theorem.
The trace assertion follows from its entrywise residual
triple.  The asserted $K$-group vanishing and the certified
primality are KT-ORDINARY-SYM2-EXPLICIT-LARGE.
Finally \texttt{CoversPt} covers every \texttt{TwoPt}
up to torus conjugation, which preserves the isomorphism
class on $G_E$.
\end{proof}

\paragraph{What the sign distinction does and does not prove.}
The actual source points above rule out an exclusion based
only on the opposite signs in (HS.2).  The archive injects
the rational strict Selmer group into a larger weight
$-1$ deformation module and bounds that module using all
relevant points.  The presence of a point is not the
converse of a nonzero rational strict Selmer class.
A sharper argument would have to retain the first
eigensummand throughout the deformation-length comparison.
Neither a projector on its tangent space alone nor removal
of adjoint characteristic-zero points has been proved to
give such a comparison here.  Also, the existential
\texttt{OneCharU\_s} statement may admit finite sets other
than the particular \texttt{CoversPt} set constructed in
its proof; the corollary makes no assertion about every
possible existential witness.

\paragraph{Exact source and proof-check scope.}
All Lean paths below are relative to the pinned
\texttt{AutomorphicVarieties/} directory.
\texttt{UnitaryThreeGroups/UnitaryThreeReflectionRight.lean}
defines $J$, $\theta$, $\theta_{\rm Lie}$, and $r_0$.
\texttt{ShaInjGen.lean}, lines 101--129, defines and checks
the Selmer matrix in (HS.2).
\texttt{BaseRepGL3NormalisedOfPsiEqOne.lean}, lines 15--24,
defines \texttt{TypeI}, \texttt{Upper}, and \texttt{Lower}.
\texttt{NormEx.lean}, lines 155 onward, proves universal
normalizability by a matrix congruent to one.
\texttt{GRepresents.lean}, lines 37--105, and
\texttt{ConjStable.lean}, lines 49 and 94--99, give the
universal ring and the literal point-cover definitions.
\texttt{RhoKApply.lean}, lines 61--83, defines
$D.\mathrm{Kill}$ and $D.\mathrm{IsPt}$; the one-way
\texttt{isPt\_push} proof is in \texttt{IsPtOf.lean},
lines 124--166.

The local ideals are defined in
\texttt{PadicHodgeTheory/IsInertiaInfiniteCyclotomic.lean}.
The complete proof of (HS.5) is the carried
\texttt{component\_basic} theorem in
\texttt{ModularityArgument/CrystallineOrdinaryIdeals/}
\texttt{OrdinaryComponentUniquePrimeDimThirteenNearBase.lean};
its local published inputs remain explicit dependencies.
The saturation calculation is checked against that
directory's \texttt{SaturatingElementUnramifiedTwistEqPowMulUnit.lean}.
The seemingly relevant \texttt{OnOrdinaryComponentIffContinuousCrystallineLift}
theorem is restricted to Artinian characteristic-zero
deformations of the exact base representation and is not
used as an arbitrary ordinary-point converse.
The separately named \texttt{UnitaryThreeDeformations/IsHostFrame}
fixes the Gaussian character and is not substituted for
the present Eisenstein-field data.

Patrikis's Proposition~4.4 and its entire proof were read
in the primary preprint, arXiv:1507.01294, pp.~12--13;
the elementary extension of its obstruction proof to
$B\rtimes\eta$ is written out above.  The original local
ordinary theorem for modular forms and the all-prime
non-CM image theorem remain the precisely named primary
inputs already recorded in KT-ORDINARY-SYM2-FAMILY.
The separate \texttt{host\_sign\_matrix\_check.py/json}
certificate verifies sixteen matrix and scalar identities
over the rational function field with $2s^2+1=0$, including
both $\eta$ components of (HS.14), the residual cancellation,
the root signs, and (HS.10).  This finite algebra check is
not presented as a verification of the deformation-theoretic
or automorphic inputs.
